\documentclass[preprint,12pt]{elsarticle}

\usepackage{booktabs}
\usepackage{url}
\usepackage[hidelinks]{hyperref}
\usepackage{xcolor}

% Un hueco marcado es un hueco. Una estimación con aire de medición es un error.
\newcommand{\pendiente}[1]{\textcolor{red}{[\textbf{pendiente:} #1]}}

\journal{Journal of Systems and Software}

\begin{document}

\begin{frontmatter}

\title{Removing the servlet container: architecture and a verification
       method for a dependency-free Java web framework}

\author[una]{Richar Andre Vilca-Solorzano\corref{cor}}
\ead{75521963@est.unap.edu.pe}
\ead[orcid]{0009-0003-2385-5263}
\cortext[cor]{Corresponding author.}

\author[una]{Fred Torres-Cruz}
\ead[orcid]{0000-0003-0834-6834}

\author[una]{Ramiro Pedro Laura-Murillo}
\ead[orcid]{0000-0003-1837-4871}

\affiliation[una]{organization={Universidad Nacional del Altiplano,
                                Facultad de Ingeniería Estadística e Informática},
                  city={Puno}, country={Peru}}

\begin{abstract}
\pendiente{escribir al final, cuando los resultados estén cerrados}
\end{abstract}

\begin{keyword}
web framework \sep servlet container \sep Java \sep virtual threads \sep
software architecture \sep software verification
\end{keyword}

\end{frontmatter}

\section{Introduction}
\label{sec:introduction}

Deploying a Java web application has long implied a servlet container: package a
WAR, install it into Tomcat or an equivalent, and rely on server configuration
that lives outside the program. None of that is the program the developer wrote.

This paper reports the design of \emph{Cero}, a Java web framework with its own
HTTP server and no runtime dependencies beyond the JDK, and --- the part we
argue is transferable --- the method used to verify that removing the container
did not silently drop behaviour that the container used to provide.

The contribution is not the measured numbers. It is a reproducible procedure for
this class of migration, and the evidence that the procedure finds defects that
conventional test-suite growth does not.

\paragraph{Contributions}
\begin{enumerate}
  \item An architecture in which the dependency on the servlet API is reduced to
        five interfaces, with a servlet adapter retained deliberately as an exit
        door so that the migration stays reversible (Section~\ref{sec:architecture}).
  \item A verification method organised around \emph{uncovered exit paths}
        rather than uncovered lines, together with the defects it found that a
        line-coverage criterion did not (Section~\ref{sec:method}).
  \item Three checks that reveal when a comparative load benchmark is measuring
        its own harness rather than the systems under test, and the factor-of-four
        correction they produced here (Section~\ref{sec:harness}).
  \item A conformance-vector suite for RFC~9112 and RFC~9110 written from the
        specification text, and the four violations it exposed in an
        implementation that already passed its own test suite
        (Section~\ref{sec:conformance}).
  \item An evaluation against eight other JVM frameworks under identical
        conditions, reported with the measurement decisions that change its
        conclusions (Section~\ref{sec:evaluation}).
\end{enumerate}

\section{Background and related work}
\label{sec:background}

\pendiente{revisión de literatura: frameworks JVM ligeros, coste del contenedor
de servlets, hilos virtuales (JEP 444), y trabajos previos sobre conformidad de
servidores HTTP}

\section{Architecture}
\label{sec:architecture}

\subsection{What the container actually provided}

An application that appears bound to a servlet container is, in our case, bound
to five interfaces: \texttt{HttpServletRequest}, \texttt{HttpServletResponse},
\texttt{HttpSession}, \texttt{Cookie} and \texttt{Part}. Writing those in-house
is what makes the container removable.

\subsection{Module structure}

\begin{table}[ht]
\centering
\caption{Modules, with measured size and the checks that exercise each one.}
\label{tab:modules}
\begin{tabular}{lrrl}
\toprule
Module & Classes & Lines & Responsibility \\
\midrule
\texttt{cero-http}    &  43 &  6\,752 & HTTP/1.1 and HTTP/2 server \\
\texttt{cero-core}    &  71 &  7\,675 & routing, injection, JSON, security \\
\texttt{cero-data}    &  16 &  1\,937 & JDBC access, migrations, sessions \\
\texttt{cero-view}    &   9 &  1\,001 & template engine \\
\texttt{cero-launcher} &  6 &     881 & CLI, project generator, packaging \\
\texttt{cero-adapter-servlet} & 4 & 603 & the exit door \\
\texttt{cero-test}    &   4 &     439 & end-to-end support for consumers \\
\midrule
Total                 & 153 & 19\,288 & \\
\bottomrule
\end{tabular}
\end{table}

\subsection{The exit door}

\texttt{cero-adapter-servlet} lets the same application deploy inside Tomcat
without source changes. It is the only place in the project that references
\texttt{jakarta.*}, and it is declared \emph{provided}. Its purpose is not
interoperability for its own sake: it makes the decision to leave the container
reversible, which is what allows the migration to be attempted at all.

\section{Verification method}
\label{sec:method}

\subsection{Uncovered exit paths, not uncovered lines}

The two most serious defects found in this project share one shape. Both were
paths out of the system that no test traversed, not lines that no test executed.

\paragraph{The session cookie over HTTP/2}
Two sibling implementations of the same response interface existed; only one
asked for the pending session cookie. Over HTTP/2 no client could therefore
establish a session, and without a session there is no CSRF token, so every
write returned 403 pointing at the wrong cause. None of the 432 checks in the
module observed it, because all of them entered over HTTP/1.1.

\paragraph{Body validation after the response}
Finishing a response closes the local half of an HTTP/2 stream, not the stream
(RFC~9113 §5.1). The server removed the stream as soon as the handler returned,
so a \texttt{DATA} frame arriving afterwards was no longer checked against the
declared \texttt{content-length} --- the precondition for request smuggling
behind a proxy. The defect was intermittent: it appeared only when the response
won the race against the request body, which is why the module's own suite never
saw it while an external conformance suite failed roughly one run in five.

Line coverage does not distinguish either case. Counting exit paths does.

\subsection{Validating the measurement harness before trusting it}
\label{sec:harness}

A comparative benchmark is an instrument, and an instrument that has not been
checked measures itself. Ours did, in three independent ways, and all three were
invisible in the output: every run completed, no framework errored, and the
numbers looked plausible.

\paragraph{Symptom 1: added load reduces throughput and server utilisation}
Raising the connection count from 64 to 1024 lowered throughput from 76 to
61~kreq/s \emph{and} lowered the server's CPU utilisation from 74\% to 55\%.
A server that is the bottleneck saturates as load grows. One whose utilisation
\emph{falls} as load grows is not the bottleneck: the constraint lies on the
path to it.

\paragraph{Symptom 2: throughput equals concurrency divided by latency}
Measured throughput was 73\,679~req/s at a mean latency of 0.868~ms with 64
connections; $64/0.000868 = 73\,733$. When Little's law holds to four
significant figures, the experiment is reporting the offered concurrency, not
the system's capacity.

\paragraph{Symptom 3: neither endpoint is saturated}
Per-core sampling showed the load generator pinned to two cores at 98--99\%
while the server sat at 40\%. The harness was measuring its own client.

\paragraph{What the checks found}
Correcting the three yielded, for the same framework on the same machine:
client cores raised from 2 to 12 (22 to 74~kreq/s), and the container's
published port bypassed in favour of its address on the bridge network, since
the port-forwarding path capped throughput at 76~kreq/s against 94~kreq/s
direct --- 23\% consumed by the container runtime rather than by any framework
under test.

We report this because the corrected figures differ from the uncorrected ones by
a factor of four, and because the uncorrected run produced a coherent, publishable
looking table in which all nine contenders fell within 2\% of one another. That
table was an artefact. Any comparative benchmark should report the utilisation
of both endpoints and the relationship between concurrency, latency and
throughput, so that a reader can see which component was actually measured.

\subsection{Deploying for real}

Two further defects appeared only behind a reverse proxy: a session cookie
issued without \texttt{Secure} because the application did not know TLS was
terminated upstream, and a deployment script that erased the certbot HTTPS block
when rewriting the vhost.

\section{Conformance vectors}
\label{sec:conformance}

No free conformance suite exists for HTTP/1.1 servers. We wrote 23 vectors
covering 13 distinct sections of RFC~9112 and RFC~9110, each sent as raw bytes
and each citing the clause it enforces. They exposed four violations in an
implementation that already passed its own tests: \emph{absolute-form} and
\emph{asterisk-form} request targets were rejected, \texttt{chunked} was
accepted outside the final position of \texttt{Transfer-Encoding}, and null
bytes were admitted in header values.

For HTTP/2 we use h2spec, an existing suite: 146 tests, of which
\pendiente{cifra final} pass.

\section{Evaluation}
\label{sec:evaluation}

\pendiente{tabla del banco; los absolutos esperan la corrida en Linux sin
virtualizar, y la columna de peticiones por segundo no discrimina con el cliente
de carga en el mismo equipo}

\section{Threats to validity}
\label{sec:threats}

\paragraph{Measurement environment}
All comparative figures were obtained under Docker Desktop on macOS, that is,
inside a virtual machine, with the load generator on the same host. Conditions
were identical for all nine frameworks, so relative comparisons hold; absolute
figures are not citable and are reported as such.

\paragraph{A measurement decision that inverts the conclusion}
Without an explicit heap cap, the JVM sizes its maximum heap from the container
limit, and a framework that allocates more per request fills it while one that
allocates little does not. The resident-set column then compares garbage
collector behaviour rather than frameworks. Under that setting our framework
measured lowest of the group; with an identical \texttt{-Xmx} for every
contender it measures highest. We report the capped configuration and state the
cap, because a benchmark whose conclusion reverses when one variable is fixed
must be published with that variable fixed.

\paragraph{Single-team implementation}
The framework, its test suite and its conformance vectors were written by the
same authors. The external audit that produced eleven findings is the only
independent evidence reported here.

\section{Conclusion}
\label{sec:conclusion}

\pendiente{cerrar cuando estén los resultados}

\section*{CRediT authorship contribution statement}

\textbf{Richar Andre Vilca-Solorzano:} Conceptualization, Methodology,
Software, Validation, Investigation, Data curation, Writing --- original draft,
Visualization. \textbf{Fred Torres-Cruz:} Conceptualization, Methodology,
Writing --- review \& editing, Supervision. \textbf{Ramiro Pedro
Laura-Murillo:} Conceptualization, Writing --- review \& editing, Supervision.

\section*{Declaration of competing interest}

The authors declare no competing financial interests or personal relationships
that could have appeared to influence the work reported in this paper.

\section*{Note on prior work}

The first author previously contributed to JxMVC, a servlet-based MVC framework
authored by R.~P.~Laura-Murillo. Cero is a separate system with its own HTTP
server: no class inherited from that framework remains, and the measurement
establishing this is recorded in the project history.

\section*{Data availability}

Source code, test suite, conformance vectors, benchmark harness and raw
measurements are public at \url{https://github.com/Andre031222/Cero}.

\bibliographystyle{elsarticle-num}
\bibliography{referencias}

\end{document}
